%% ma: L12-B01-0004
%% lop: 12
%% chuong: 1
%% bai: 1
%% dang: 12.1.4
%% loai: TN4
%% muc: NB
%% dap_an: "D"
%% nguon: "(NBV)-ĐỀ SỐ 2-MỖI NGÀY 1 ĐỀ THI 2025.pdf (D02, MỖI NGÀY 1 ĐỀ - Đề số 2), Phần I Câu 1"
%% kien_thuc: "Bài 1 (điểm cực trị của đồ thị hàm số, đọc từ đồ thị)"
%% ghi_chu: "Hình dựng lại bằng hàm bậc ba y=x^3-3x-2 (qua (-1;0), (0;-2), (1;-4), (2;0)), khớp số liệu hình gốc"
%% trang_thai: chinh-thuc
\begin{ex}
Cho hàm số $y=f(x)$ có đồ thị như hình vẽ. Điểm cực tiểu của đồ thị hàm số đã cho là
\begin{center}
\begin{tikzpicture}[x=1cm,y=.6cm]
  \draw[truc] (-2.1,0)--(2.8,0) node[right]{$x$};
  \draw[truc] (0,-4.9)--(0,3.1) node[above]{$y$};
  \draw[smooth,samples=80,domain=-1.78:2.28] plot(\x,{(\x)^3-3*\x-2});
  \draw[phu] (1,0)--(1,-4)--(0,-4);
  \draw[phu] (-1,0)--(-1,0);
  \node[below left] at (0,0) {$O$};
  \node[above] at (-1,0) {$-1$};
  \node[above right] at (1,0) {$1$};
  \node[left] at (0,-2) {$-2$};
  \node[left] at (0,-4) {$-4$};
  \fill (0,-2) circle (1.2pt);
\end{tikzpicture}
\end{center}
\choice{$-1$}{$1$}{$(2;0)$}{\True $(1;-4)$}
\loigiai{Đồ thị chuyển từ đi lên sang đi xuống tại $x=-1$ (điểm cực đại) và từ đi xuống sang đi lên tại $x=1$. Vậy điểm cực tiểu của đồ thị là $(1;-4)$.}
\end{ex}
